% SPDX-License-Identifier: Apache-2.0
\section{The device under test}
\label{sec:dut}

The device under test is Vreteno's netlist, the Verilog that TxHDL
writes for the core. It is the same netlist that the board's
bitstream is built from, so the tests check what is synthesized rather
than a model of it.

\subsection{The wrapper}

The core on its own speaks TxHDL's transaction channels, not a standard
bus. A small TxHDL unit, \code{Dut}, joins three existing parts so that
the netlist has one standard AXI4 manager port (Fig.~\ref{fig:dut}).
The three parts are the hart, which is the core with its memory
management unit; the tracker, which turns the core's loads, stores and
fetches into AXI4 bursts; and an adapter that puts the five AXI4
channels on pins. All three come from TxHDL unchanged. The wrapper is
about 250 lines of Rust, most of them port names, and TxHDL lowers it
with the core into one Verilog file.

\begin{figure}[t]
\centering
\begin{tikzpicture}[node distance=5mm and 4mm, font=\footnotesize]
\node[box, minimum width=22mm] (hart) {hart\\\scriptsize core, MMU, I-cache};
\node[box, right=of hart, minimum width=14mm] (trk) {tracker};
\node[box, right=of trk, minimum width=14mm] (pins) {AXI4\\pins};
\node[draw, dashed, rounded corners=2pt, inner sep=3mm,
      fit=(hart)(trk)(pins), label={[lbl]above:\code{Dut}, the netlist}] (dut) {};
\node[tool, below=13mm of trk, minimum width=50mm] (mem)
  {memory model in C++\\\scriptsize ELF loader, \code{tohost}, console};
\draw[fl] (hart) -- (trk);
\draw[fl] (trk) -- (pins);
\draw[fl, <->] (pins.south) -- node[lbl, right, pos=0.72] {AXI4} (pins.south |- mem.north);
\end{tikzpicture}
\caption{The device under test. Only the memory model is written for
this testbench; everything inside the dashed box is TxHDL's.}
\label{fig:dut}
\end{figure}

\subsection{Booting into the test}

Vreteno fetches its first 4\,KiB from a memory inside the core, which
the netlist initializes. The wrapper puts two instructions there:
\code{lui t0,0x40000} and \code{jr t0}. They jump to
\code{0x4000\_0000}, where every test is linked. On the board, that
address is the start of DDR3, and the core's instruction cache holds
words from that range. So the tests run from cached memory, as the
board's programs do.

\subsection{The memory and the end of a run}

The testbench is a C++ program around the Verilated netlist. It loads
the test's ELF into a sparse memory and answers every AXI4 burst from
it, with every \code{ready} held high. The interrupt lines and the
debug module's lines are held low.

A test ends the way it ends on Sail, through the HTIF \code{tohost}
word. The testbench reads the word's address from the ELF's symbol
table. A write to its upper half completes a command. Device~1 with
command~1 prints a character on the console. Any other command with an
odd lower half ends the run: 1 means the test passed, and 3 means it
failed.

The testbench writes a result file for each run: the status, the
cycles, the instructions retired, and the console text. A run that
reaches 20 million cycles stops with the status \code{TIMEOUT}. The
build does not stop when a test fails. The result file records the
failure, a test target reads it, and this report counts it.
